% =====================================================================
%  Tarea #1 - Vibracion libre de 1 GDL, sistemas discretos
%  Vibraciones Mecanicas - Universidad Rafael Landivar
%  Santiago Cruz Rivera - Carne 1310823
%  Compilar con: pdflatex Tarea1_VM.tex   (dos veces)
% =====================================================================
\documentclass[12pt,letterpaper]{article}

\usepackage[utf8]{inputenc}
\usepackage[T1]{fontenc}
\usepackage[spanish,es-tabla]{babel}
\usepackage[margin=2.5cm]{geometry}
\usepackage{amsmath,amssymb}
\usepackage{graphicx}
\usepackage{xcolor}
\usepackage{listings}
\usepackage{float}
\usepackage{caption}

% ------------------------- Colores y estilo de codigo ----------------
\definecolor{mlkeyword}{RGB}{0,0,255}
\definecolor{mlcomment}{RGB}{34,139,34}
\definecolor{mlstring}{RGB}{160,32,240}
\definecolor{fondo}{RGB}{248,248,248}
\definecolor{urlazul}{RGB}{0,51,102}

\lstdefinestyle{matlab}{
  language=Matlab,
  backgroundcolor=\color{fondo},
  basicstyle=\ttfamily\footnotesize,
  keywordstyle=\color{mlkeyword},
  commentstyle=\color{mlcomment}\itshape,
  stringstyle=\color{mlstring},
  numbers=left,
  numberstyle=\tiny\color{gray},
  numbersep=8pt,
  frame=single,
  rulecolor=\color{gray},
  breaklines=true,
  breakatwhitespace=false,
  showstringspaces=false,
  tabsize=2,
  captionpos=b,
  columns=flexible,
  literate={á}{{\'a}}1 {é}{{\'e}}1 {í}{{\'i}}1 {ó}{{\'o}}1 {ú}{{\'u}}1 {ñ}{{\~n}}1
}
\lstset{style=matlab}

% ------------------------- Encabezado de secciones -------------------
\usepackage{titlesec}
\titleformat{\section}{\large\bfseries\color{urlazul}}{Problema \thesection.}{0.5em}{}
\titleformat{\subsection}{\normalsize\bfseries}{}{0em}{}

\usepackage{fancyhdr}
\pagestyle{fancy}
\fancyhf{}
\fancyhead[L]{\footnotesize Vibraciones Mecánicas --- Tarea \#1}
\fancyhead[R]{\footnotesize Santiago Cruz Rivera --- 1310823}
\fancyfoot[C]{\thepage}
\renewcommand{\headrulewidth}{0.4pt}

\begin{document}

% =====================================================================
%                              CARÁTULA
% =====================================================================
\begin{titlepage}
\centering
\thispagestyle{empty}

\vspace*{1.5cm}

{\LARGE\bfseries UNIVERSIDAD RAFAEL LANDÍVAR\par}
\vspace{0.4cm}
{\large Facultad de Ingeniería\par}

\vspace{2.5cm}

\rule{\textwidth}{1.2pt}
\vspace{0.6cm}

{\Huge\bfseries VIBRACIONES MECÁNICAS\par}
\vspace{0.8cm}
{\LARGE TAREA \#1\par}
\vspace{0.5cm}
{\Large Vibración libre de 1 GDL --- Sistemas discretos\par}

\vspace{0.6cm}
\rule{\textwidth}{1.2pt}

\vspace{3cm}

{\Large\bfseries Santiago Cruz Rivera\par}
\vspace{0.4cm}
{\large Carné: 1310823\par}

\vfill

{\large Guatemala, \today\par}

\end{titlepage}

% =====================================================================
%                            PROBLEMA 1
% =====================================================================
\section{Movimiento armónico de una mesa vibratoria}

\subsection{Enunciado}
\begin{figure}[H]
\centering
\includegraphics[width=0.95\textwidth]{img/ej1.png}
\caption*{Enunciado del problema 1.}
\end{figure}

\subsection{Fórmulas utilizadas}

Frecuencia angular a partir de la frecuencia cíclica:
\begin{equation}
\omega = 2\pi f
\end{equation}

Para que el paquete $B$ no deslice sobre la mesa $C$, la fuerza de fricción estática
debe ser capaz de proporcionar la aceleración del paquete. En el límite de deslizamiento:
\begin{equation}
F_{\max} = m\,a_{\max}, \qquad N = mg, \qquad F_{\max} = \mu_s N
\end{equation}

Sustituyendo, la masa se cancela y la aceleración máxima admisible resulta:
\begin{equation}
m\,a_{\max} = \mu_s\,m\,g \quad \Longrightarrow \quad a_{\max} = \mu_s\, g
\end{equation}

En movimiento armónico simple, $x(t) = A\sin(\omega t)$, por lo que la aceleración máxima es:
\begin{equation}
a_{\max} = \omega^{2} A
\end{equation}

Igualando ambas expresiones se despeja la amplitud máxima permitida:
\begin{equation}
\boxed{\;A_{\max} = \frac{\mu_s\, g}{\omega^{2}} = \frac{\mu_s\, g}{(2\pi f)^{2}}\;}
\end{equation}

Conversión al sistema inglés:
\begin{equation}
A_{\text{US}} \,[\text{in}] = A\,[\text{m}] \times 39.37
\end{equation}

\subsection{Código MATLAB}
\lstinputlisting[firstline=1,lastline=12,firstnumber=1]{Tarea1_VM.m}

\subsection{Resultados}
\begin{figure}[H]
\centering
\includegraphics[width=0.62\textwidth]{R1.png}
\caption*{Salida de la ventana de comandos --- Problema 1.}
\end{figure}

\noindent Con $f = 3$ Hz, $\mu_s = 0.40$ y $g = 9.81$ m/s$^2$:
\[
\omega = 18.850\ \text{rad/s}, \qquad a_{\max} = 3.924\ \text{m/s}^2,
\qquad A_{\max} = 11.04\ \text{mm} = 0.4348\ \text{in}
\]

\newpage
% =====================================================================
%                            PROBLEMA 2
% =====================================================================
\section{Bloque de 75 lb con arreglos de resortes}

\subsection{Enunciado}
\begin{figure}[H]
\centering
\includegraphics[width=0.85\textwidth]{img/ej2.png}
\caption*{Enunciado del problema 2.}
\end{figure}

\subsection{Fórmulas utilizadas}

\textbf{Masa a partir del peso} (sistema inglés, $g = 386.4$ in/s$^2$):
\begin{equation}
m = \frac{W}{g} = \frac{75}{386.4} = 0.19410\ \frac{\text{lb}\cdot\text{s}^2}{\text{in}}
\end{equation}

\textbf{Rigidez equivalente.} Resortes en paralelo (mismo desplazamiento) y en serie
(misma fuerza):
\begin{equation}
k_{eq}^{\;\text{paralelo}} = \sum_{i} k_i, \qquad
\frac{1}{k_{eq}^{\;\text{serie}}} = \sum_{i}\frac{1}{k_i}
\end{equation}

\emph{Sistema 1:} el resorte superior de 90 lb/in y los dos resortes inferiores de
45 lb/in actúan todos en paralelo sobre el bloque:
\begin{equation}
k_{e1} = 90 + 2(45) = 180\ \text{lb/in}
\end{equation}

\emph{Sistema 2:} los dos resortes de 90 lb/in están en serie:
\begin{equation}
k_{e2} = \frac{1}{\dfrac{1}{90}+\dfrac{1}{90}} = \frac{90}{2} = 45\ \text{lb/in}
\end{equation}

\textbf{Parámetros del movimiento armónico simple:}
\begin{align}
\omega_n &= \sqrt{\frac{k_{eq}}{m}} \\[4pt]
f &= \frac{\omega_n}{2\pi} \\[4pt]
T &= \frac{1}{f} = \frac{2\pi}{\omega_n} \\[4pt]
v_{\max} &= A\,\omega_n \\[4pt]
a_{\max} &= A\,\omega_n^{2}
\end{align}

\subsection{Código MATLAB}
\lstinputlisting[firstline=14,lastline=40,firstnumber=14]{Tarea1_VM.m}

\subsection{Resultados}
\begin{figure}[H]
\centering
\includegraphics[width=\textwidth]{R2.png}
\caption*{Salida de la ventana de comandos --- Problema 2.}
\end{figure}

\begin{center}
\begin{tabular}{lcc}
\hline
\textbf{Parámetro} & \textbf{Sistema 1} & \textbf{Sistema 2} \\
\hline
$k_{eq}$ [lb/in]        & 180    & 45     \\
$\omega_n$ [rad/s]      & 30.45  & 15.23  \\
$f$ [Hz]                & 4.85   & 2.42   \\
$T$ [s]                 & 0.21   & 0.41   \\
$v_{\max}$ [in/s]       & 60.91  & 30.45  \\
$a_{\max}$ [in/s$^2$]   & 1854.72& 463.68 \\
\hline
\end{tabular}
\end{center}

\newpage
% =====================================================================
%                            PROBLEMA 3
% =====================================================================
\section{Determinación de la masa del bloque \emph{C} y su período}

\subsection{Enunciado}
\begin{figure}[H]
\centering
\includegraphics[width=0.95\textwidth]{img/ej3.png}
\caption*{Enunciado del problema 3.}
\end{figure}

\subsection{Fórmulas utilizadas}

Relación entre período y frecuencia natural:
\begin{equation}
T = \frac{2\pi}{\omega_n} \quad \Longrightarrow \quad \omega_n = \frac{2\pi}{T}
\end{equation}

Frecuencia natural del sistema masa--resorte:
\begin{equation}
\omega_n = \sqrt{\frac{k}{m}} \quad \Longrightarrow \quad k = \omega_n^{2}\, m
\end{equation}

Como la rigidez $k$ del arreglo de resortes no cambia al retirar bloques, se plantea
un sistema de dos ecuaciones con dos incógnitas ($k$ y $m_C$):

\emph{Condición 1} --- con los bloques $A$, $B$ y $C$ ($T_1 = 0.8$ s):
\begin{equation}
k = \omega_{1}^{2}\,(m_A + m_B + m_C), \qquad \omega_1 = \frac{2\pi}{0.8}
\end{equation}

\emph{Condición 2} --- retirando el bloque $A$ ($T_2 = 0.7$ s):
\begin{equation}
k = \omega_{2}^{2}\,(m_B + m_C), \qquad \omega_2 = \frac{2\pi}{0.7}
\end{equation}

Igualando ambas expresiones de $k$:
\begin{equation}
\omega_1^{2}(m_A+m_B+m_C) = \omega_2^{2}(m_B+m_C)
\quad\Longrightarrow\quad
\boxed{\;m_C = \frac{\omega_1^{2}\,m_A - (\omega_2^{2}-\omega_1^{2})\,m_B}{\omega_2^{2}-\omega_1^{2}}\;}
\end{equation}

\emph{Condición 3} --- retirando $A$ y $B$, queda únicamente $C$:
\begin{equation}
\omega_3 = \sqrt{\frac{k}{m_C}}, \qquad T_3 = \frac{2\pi}{\omega_3}
\end{equation}

En MATLAB el sistema se resuelve simbólicamente con \texttt{syms} y \texttt{solve}.

\subsection{Código MATLAB}
\lstinputlisting[firstline=42,lastline=64,firstnumber=42]{Tarea1_VM.m}

\subsection{Resultados}
\begin{figure}[H]
\centering
\includegraphics[width=0.52\textwidth]{r3.png}
\caption*{Salida de la ventana de comandos --- Problema 3.}
\end{figure}

\noindent Con $\omega_1 = 7.854$ rad/s y $\omega_2 = 8.976$ rad/s:
\[
k = 789.57\ \text{N/m}, \qquad m_C = 6.80\ \text{kg}, \qquad
\omega_3 = 10.77\ \text{rad/s}, \qquad T_3 = 0.58\ \text{s}
\]

\newpage
% =====================================================================
%                            PROBLEMA 4
% =====================================================================
\section{Viga en voladizo como resorte equivalente}

\subsection{Enunciado}
\begin{figure}[H]
\centering
\includegraphics[width=0.95\textwidth]{img/ej4.png}
\caption*{Enunciado del problema 4.}
\end{figure}

\subsection{Fórmulas utilizadas}

\textbf{(a) Constante de resorte equivalente.} De mecánica de materiales, la deflexión
en el extremo libre $B$ de una viga en voladizo con carga puntual $P$ es:
\begin{equation}
\delta_B = \frac{P L^{3}}{3EI}
\end{equation}

Por definición de rigidez, $k = P/\delta$, por lo tanto:
\begin{equation}
\boxed{\;k_{eq} = \frac{P}{\delta_B} = \frac{3EI}{L^{3}}\;}
\end{equation}

Con $L = 10\ \text{ft} = 120\ \text{in}$, $E = 29\times10^{6}\ \text{lb/in}^2$ e
$I = 12.4\ \text{in}^4$.

\textbf{(b) Frecuencia de vibración.} La masa del bloque de 520 lb es:
\begin{equation}
m = \frac{W}{g} = \frac{520}{386.4} = 1.3457\ \frac{\text{lb}\cdot\text{s}^2}{\text{in}}
\end{equation}

\begin{equation}
\omega_n = \sqrt{\frac{k_{eq}}{m}}, \qquad
f = \frac{\omega_n}{2\pi}, \qquad
T = \frac{2\pi}{\omega_n}
\end{equation}

\subsection{Código MATLAB}
\lstinputlisting[firstline=66,lastline=75,firstnumber=66]{Tarea1_VM.m}

\subsection{Resultados}
\begin{figure}[H]
\centering
\includegraphics[width=0.52\textwidth]{r4.png}
\caption*{Salida de la ventana de comandos --- Problema 4.}
\end{figure}

\noindent Sustituyendo valores:
\[
k_{eq} = \frac{3(29\times10^{6})(12.4)}{120^{3}} = 624.31\ \text{lb/in}
\]
\[
\omega_n = \sqrt{\frac{624.31}{1.3457}} = 21.54\ \text{rad/s}
\]
\[
f = \frac{\omega_n}{2\pi} = 3.43\ \text{Hz}, \qquad
T = \frac{2\pi}{\omega_n} = 0.29\ \text{s}
\]

\end{document}
